\section{Постановка задачи}

\subsection{Цель}
Построить и проанализировать локальное двумерное сечение функции потерь обученной CNN-модели на CIFAR-10, сравнив качество трёх квадратичных аппроксимаций поверхности.

\subsection{Формальная постановка}
Рассматривается поверхность:
\[
f(\alpha, \beta) = L(\theta + \alpha d_1 + \beta d_2),
\]
где $\theta$ --- параметры обученной модели, а $d_1, d_2$ --- случайные направления в пространстве параметров после filter-wise normalization.

Задача:
\begin{enumerate}
    \item вычислить $f(\alpha,\beta)$ на сетке в диапазоне $[-r, r]^2$;
    \item получить 1D-срезы $f(\alpha,0)$ и $f(0,\beta)$;
    \item сравнить аппроксимации $q_1, q_2, q_3$ по метрикам $RMSE$, $L_\infty$, $RelRMSE$;
    \item интерпретировать, как меняется ошибка аппроксимации при изменении радиуса.
\end{enumerate}

\subsection{Гипотеза исследования}
В малой окрестности минимума полная квадратичная аппроксимация ($q_1$) и диагональная аппроксимация ($q_2$) должны описывать поверхность точнее, чем аппроксимация на базе конечных разностей гессиана ($q_3$), чувствительная к численным эффектам.
